\chapter{Conclusiones}

Esta tesis presentó un estudio de la retroproyección filtrada en reconstrucción tomográfica bidimensional, combinando derivación matemática e implementación reproducible. En el modelo continuo se fijaron las convenciones de Fourier, la parametrización normal de las rectas y el intervalo angular $\theta\in[0,\pi)$. Con ellas, la transformada de Radon se relacionó con la transformada de Fourier bidimensional mediante el teorema de corte, y la fórmula de FBP se obtuvo por el cambio a coordenadas polares con radio firmado. Así se identificó el filtro rampa $\tau_{\mathrm{rampa}}(\sigma)=|\sigma|$ como el factor espectral asociado al jacobiano radial, sin factores adicionales de normalización.

La parte computacional trasladó esta formulación a un esquema discreto reproducible con un fantoma geométrico propio, sinogramas de haz paralelo bajo la hipótesis de soporte en el círculo inscrito, filtrado por columnas mediante FFT con relleno centrado de ceros y retroproyección discreta. La rampa se mantuvo como referencia teórica y numérica, mientras que los filtros A--F se usaron como modificaciones exploratorias del símbolo para observar cómo distintas respuestas espectrales afectan la reconstrucción.

Los experimentos muestran el compromiso esperado entre resolución y estabilidad. Sin ruido, los filtros D y E obtuvieron los menores errores relativos y quedaron prácticamente empatados; la rampa produjo un resultado cercano y conservó su papel de referencia. Con ruido, bajo la configuración y la realización aleatoria consideradas, el filtro C alcanzó el menor error relativo. Esta observación no implica superioridad universal del filtro C, sino una conclusión restringida al fantoma, al nivel de ruido, a la semilla y a los parámetros numéricos utilizados.

Las limitaciones del estudio son deliberadas: geometría paralela bidimensional, fantoma sintético, una sola realización de ruido y filtros exploratorios. No se usaron datos clínicos ni tridimensionales, no se ajustó empíricamente la amplitud de las reconstrucciones y no se intentó demostrar optimalidad. Estas restricciones permitieron mantener una conexión clara entre la fórmula continua de FBP, su implementación discreta y la interpretación de los resultados.

Como trabajo futuro, sería natural ampliar los experimentos a múltiples semillas aleatorias, otros niveles de ruido, otras familias de fantomas y criterios cuantitativos adicionales. También podría estudiarse el diseño sistemático de filtros bajo objetivos de error o estabilidad, así como configuraciones con datos incompletos o de ángulo limitado \cite{quinto1993,frikel_quinto2013}.
